\begin{lemma}
\label{lemma-in-image}
Let $\varphi : R \to S$ be a ring map. Let $\mathfrak p$
be a prime of $R$. The following are equivalent
\begin{enumerate}
\item $\mathfrak p$ is in the image of
$\Spec(S) \to \Spec(R)$,
\item $S \otimes_R \kappa(\mathfrak p) \not = 0$,
\item $S_{\mathfrak p}/\mathfrak p S_{\mathfrak p} \not = 0$,
\item $(S/\mathfrak pS)_{\mathfrak p} \not = 0$, and
\item $\mathfrak p = \varphi^{-1}(\mathfrak pS)$.
\end{enumerate}
\end{lemma}

\begin{proof}
We have already seen the equivalence of the first two
in Remark \ref{remark-fundamental-diagram}.
Put $T=R\setminus\mathfrak p$ and $Q=S/\mathfrak pS$, with
$T$ acting on $Q$ through $\varphi$.
The identifications in that remark identify the three rings in
(2), (3), and (4) with $T^{-1}Q$. In particular the first map is
$$
S\otimes_R\kappa(\mathfrak p)\longrightarrow T^{-1}Q,\qquad
s\otimes\overline{a/u}\longmapsto
\frac{\varphi(a)s+\mathfrak pS}{u},
$$
where $a\in R$, $u\in T$, and the bar denotes the class in
$R_{\mathfrak p}/\mathfrak pR_{\mathfrak p}$.
The formula is $R$-balanced, respects multiplication, kills
$\mathfrak pR_{\mathfrak p}$, and inverts every member of $T$.
Its inverse sends $(s+\mathfrak pS)/u$ to
$s\otimes\overline{1/u}$: the relations defining the quotient
and localization hold because $\mathfrak p$ acts as zero and
each $u\in T$ as a unit in the second tensor factor.
Both composites are the identity on the displayed generators.
The corresponding quotient-localization map is
$s/\varphi(u)+\mathfrak pS_{\mathfrak p}\mapsto
(s+\mathfrak pS)/u$, with the reverse formula as inverse.
Finally, $T^{-1}Q$ is the zero ring if and only if there is
$u\in T$ with $\varphi(u)\in\mathfrak pS$, by the equality
criterion for $1/1=0/1$ in a localization.
The ideal $\varphi^{-1}(\mathfrak pS)$ always contains
$\mathfrak p$. Hence it misses $T$ if and only if it equals
$\mathfrak p$. This proves the equivalence with (5).
\end{proof}